\section{Quantum Rejection Sampling}\label{sec:sampling}

We present the Quantum Rejection Sampling method in two steps:
\begin{itemize}
    \item {\bf Ancilla augmentation:} Find an ancilla augmentation of a distribution, such that there is a unitary $\unitarysymbol$ preparing a state with the augmented distribution as measurement distribution
    \item {\bf Acceptance amplification:} Amplify the amplitude of the ancilla qubits in state $1$ by repeated application of the Grover operator
\end{itemize}
Compared with the rejection sampling approach on the start augmentation of the distribution, we achieve a square-root speed-up in the expected number of gates to be applied per sample.

We in particular apply this scheme to sample from \ComputationActivationNetworks{}, where we apply the q-sample preparation scheme of \secref{sec:canQSample}.

\subsection{Acceptance Amplification}

Let there be a quantum state, which measurement distribution is an ancilla augmentation $\secprobat{\avariables,\shortcatvariables}$ of a target distribution $\probwith$.
We now show how Grover operators can be used to amplify the acceptance state and improve the efficiency of the rejection sampling method through an increase of the acceptance probability.
This is an instance of amplitude amplification.

\subsubsection{Grover construction}

We relate the state $\qstateof{\goodstatesymbol}$ with the event $\avariables=\onetuple{[\seldim]}$, by constructing it as
\begin{align*}
    \qstateofat{\goodstatesymbol}{\avariables,\headvariables,\shortcatvariables}
    =\onehotmapofat{\onetuple{[\seldim]}}{\avariables}
    \otimes \contractionof{\bencodingofat{\sstat}{\headvariables,\shortcatvariables},\sqrt{\probwith}}{\headvariables,\shortcatvariables}
\end{align*}
The complement state $\qstateof{\badstatesymbol}$ is constructed if $\acceptanceprob\neq1$ (in which case the acceptance amplilification scheme would not be necessary)
\begin{align*}
    \probofat{\badstatesymbol}{\avariables,\headvariables,\shortcatvariables}
    = \frac{1}{1-\acceptanceprob}
    \cdot\left(
             \secprobat{\avariables,\headvariables,\shortcatvariables}-\acceptanceprob \cdot \onehotmapofat{\onetuple{[\seldim]}}{\avariables}\otimes\probat{\headvariables,\shortcatvariables}
    \right)
\end{align*}
and
\begin{align*}
    \qstateofat{\badstatesymbol}{\avariables,\headvariables,\shortcatvariables}=\sqrt{\probofat{\perp}{\avariables,\headvariables,\shortcatvariables}} \, .
\end{align*}

% Grover operator instance
%We build the Grover operator (see \secref{sec:groverOperator}) with respect to
%\begin{itemize}
%%    \item $$
%    \item $\qstateof{\badstatesymbol}=\sqrt{\probof{\perp}}$
%\end{itemize}
%and the state

We thus have
\begin{align*}
    \qstateofat{0}{\avariables,\headvariables,\shortcatvariables}
    = \sinof{\frac{\rotanglesymbol}{2}} \qstateofat{0}{\avariables,\headvariables,\shortcatvariables}
    + \cosof{\frac{\rotanglesymbol}{2}} \qstateofat{1}{\avariables,\headvariables,\shortcatvariables}
\end{align*}
where the $\rotanglesymbol$ is related to the acceptance probability as
\begin{align*}
    \sinof{\frac{\rotanglesymbol}{2}}
    = \sqrt{\acceptanceprob} \, .
\end{align*}



\subsubsection{Circuit Realization}

Both reflection operators are realized as their negations using the Pauli-Z scheme (see \secref{sec:groverOperator}).
The negations chancel in their concatenation, which is therefore the Grover operator.


\subsubsection{Ancilla augmentation preserved}

The measurement distribution to any quantum state in $\subspaceof{}\backslash\spanof{\qstateof{\badstatesymbol}}$ is an ancilla augmentation of $\probwith$. % Need to show?
When amplifying the acceptance probability we stay in the subspace $\subspaceof{}$.
When we do not overrotate and since we start with $\acceptanceprob>0$ the amplified state is never in $\spanof{\qstateof{\badstatesymbol}}$.

%% Realization of the grover operator
%While the reflection $\reflectionof{0}$ is realized as described in \secref{sec:groverOperator}, the reflection $\reflectionof{\badstatesymbol}$ is a multi-controlled Pauli-Z on the ancilla variables. \
%\red{No, we use the negated about the good state!}



\subsection{Optimal rotation numbers}

% Application
Now the application of $\repnum$ grover operators $\groversymbol$ on $\qstateof{0}$ prepares a state
\begin{align*}
    \qstateofat{\repnum}{\avariables,\shortcatvariables}
    = \sinof{\left(\frac{1}{2}+\repnum\right)\rotanglesymbol} \qstateofat{\goodstatesymbol}{\avariables,\shortcatvariables}
    + \cosof{\left(\frac{1}{2}+\repnum\right)\rotanglesymbol} \qstateofat{\badstatesymbol}{\avariables,\shortcatvariables}
\end{align*}
The state is closest to $\qstateof{\goodstatesymbol}$ when
\begin{align*}
    \left(\frac{1}{2}+\repnum\right)
    \rotanglesymbol \approx \frac{\pi}{2} \, .
\end{align*}

Estimating for small $\acceptanceprob$ (tight for small $\acceptanceprob$)
\begin{align*}
    \sqrt{\acceptanceprob} = \sinof{\frac{\rotanglesymbol}{2}}\leq \frac{\rotanglesymbol}{2}
\end{align*}
we get a bound
\begin{align*}
    \repnumof{\mathrm{opt}} \leq \frac{\pi}{4\sqrt{\acceptanceprob}} \, .
\end{align*}
This shows that we have a square root improvement on classical rejection sampling, which would require $\frac{1}{\acceptanceprob}$ repetitions to produce a sample.
Note however that we compare quantum circuit length with the number of classical repetitions in sampling.


%Note, that the variable qubits are uniformly distributed when only the computation circuit is applied.
%When sampling the probability distribution, we need the ancilla qubits to be in state $1$ in order for the sample to be valid.
%Any other states will have to be rejected.

%\textbf{Open Question:} Is there a way to avoid amplitude amplification and use a more direct circuit implementation of the activation network?
%- Cannot be the case, when the encoding is determined by the activation tensor alone: Needs to use the computated statistic as well.
%Negative result in \cite[Section~6.6]{nielsen_quantum_2010}: When using multiple applications of \computationCircuit{} in combination with further unitaries, need at least as many applications as with amplitude amplification.

\subsection{Applications on \CompActNets{}}

The above scheme prepares the ancilla augmented \ComputationActivationNetwork{}
\begin{align*}
    \secprobat{\avariables,\headvariables,\shortcatvariables}
    = \acceptanceprob \cdot \onehotmapofat{\onetuple{[\seldim]}}{\avariables}
    \otimes \contractionof{\bencodingofat{\sstat}{\headvariables,\shortcatvariables},\probwith}{\headvariables,\shortcatvariables}
    + (1-\acceptanceprob) \probofat{\perp}{\shortcatvariables,\headvariables,\avariables}
\end{align*}
and the prepared state is the combination of q-samples
\begin{align*}
    \qstateofat{0}{\avariables,\headvariables,\shortcatvariables}
    = \sqrt{\acceptanceprob} \cdot \onehotmapofat{\onetuple{[\seldim]}}{\avariables}
    \otimes \contractionof{\bencodingofat{\sstat}{\headvariables,\shortcatvariables},\sqrt{\probwith}}{\headvariables,\shortcatvariables}
    + \sqrt{(1-\acceptanceprob)} \sqrt{\probofat{\perp}{\shortcatvariables,\headvariables,\avariables}} \, .
\end{align*}
%Here we understand $\headvariables$ as extending $\shortcatvariables$.

Sampling from the augmented \ComputationActivationNetwork{} can be also done classically by forward sampling:
Just draw the variables from uniform, calculate the value qubit by a logical circuit inference and accept with probability by the computed value.
To achieve a quantum advantage we manipulate the circuit first by the amplitude amplification scheme and increase the acceptance probability.

\subsection{Examples}

\begin{example}[Sample a model of a propositional formula]

    Acceptance rate:
    \begin{align*}
        \acceptanceprob
        = \frac{\contraction{\formula}}{2^{\catorder}}
    \end{align*}

    Angle:
    \begin{align*}
        \rotanglesymbol
        = 2 \cdot arcsin\left(\sqrt{\frac{\contraction{\formula}}{2^{\catorder}}}\right)
    \end{align*}

\end{example}

\begin{example}[Sampling from the toy accounting]

    For positive $\canparam$: $\repnumof{\mathrm{opt}}=0$

    For large negative $\canparam$: Effectively only one model (i.e. minimal acceptance prob), $\repnumof{\mathrm{opt}}=1$

\end{example}


